\scanpage{125}
\mspage{119}
% archive sheet number: 124

fin de compte, je ne suis intéressé aux supports et opérations sur les supports, que dans le cas des supports constructibles (ou ``ens. constr.'' --- \unsure{en ajoutant} encore $\unsure{\ell|\complement}\operatorname{supp}^{\unsure{0|\circ}}\unsure{\Phi|\varphi|\phi}$\ldots), et des familles comm. de tels supports.

Il y a \deleted{\unsure{donc}} là une contradiction interne, qui reste non résolue pour le moment. D'ailleurs, même la donnée \unsure{primitive} $X\unsure{\mid\circ\mid|\mid0\mid}Y$ implique la comparaison \unsure{de|des} strates $X,Y$ qui (alors $=$ qu'elles soient ``int. disj.'') peuvent ne pas être compatibles, ni même commensurables.

Pourtant, il semble indispensable d'avoir une description 'intrinsèque' des supports \unsure{subordonnés} à une figure, comme objet d'un ens. (en l'occurrence $\Sigma_{\unsure{\mathcal A|A}}$) \uline{indépendant} de ladite figure. On va \unsure{jeter|faire}
